\section*{Chapter \ref{ch:qm:markov-chains-and-queues} --- Markov Chains and Queues}

\begin{solution}{exo:qm:markov-chains-and-queues:1}
Mean holding times $1/q_i$: 0.5 s at one tick, $1/6$ s at two, $2/9 = 0.22$ s at three. From three ticks the chain goes to
two with probability $4/4.5 = 0.889$ and to one with $0.111$.
\end{solution}

\begin{solution}{exo:qm:markov-chains-and-queues:2}
$\rho = 0.8$: $L = 0.8/0.2 = 4$ messages, $W = 1/(1 - 0.8) = 5$ microseconds, and $L = \lambda W$.
\end{solution}

\begin{solution}{exo:qm:markov-chains-and-queues:3}
$(1 - (1/0.9)^{20})/(1 - (1/0.9)^{40}) = 10.8\%$.
\end{solution}

\begin{solution}{exo:qm:markov-chains-and-queues:4}
No. $\pi_1q_{12} = 0.676 \times 2 = 1.35$ but $\pi_2q_{21} = 0.265 \times 5 = 1.32$: \omterm{def:qm:markov-chains-and-queues:balance}{detailed balance} fails. Kolmogorov's criterion
compares $q_{12}q_{23}q_{31} = 2 \times 1 \times 0.5 = 1$ with $q_{13}q_{32}q_{21} = 0$: probability circulates $1 \to 2 \to 3 \to 1$.
\end{solution}

\begin{solution}{exo:qm:markov-chains-and-queues:5}
Mean $20/(1 - 0.9) = 200$ s, median 137 s, and 15.6\% within a minute. The law is skewed to the right: most paths
empty sooner than average, a few wander far up first and take very long.
\end{solution}

\begin{solution}{exo:qm:markov-chains-and-queues:6}
By Little's law $W = L/\lambda = 400/50 = 8$ seconds.
\end{solution}

\begin{solution}{exo:qm:markov-chains-and-queues:7}
73.4\% in 20\,000 seeded races, against 73.2\% by uniformisation: within the simulation's standard error of 0.3
point.
\end{solution}

\begin{solution}{exo:qm:markov-chains-and-queues:8}
The 95\% is the probability that the next tick is up, not a fill rate: an order at the back of the heavy bid fills
before the thin ask empties only 73\% of the time, and less if it is further back. And fills are not a random
sample: an order at the bid fills when sellers keep hitting it, often just before the price moves down. The
decision needs the mark-out after fills, not the fill rate.
\end{solution}

\begin{solution}{pb:qm:markov-chains-and-queues:1}
\textbf{1.} 80 orders ahead; 20 on the ask.
\textbf{2.} $80/1 + 1/0.6 = 81.7$ seconds.
\textbf{3.} Mean $20/(1 - 0.9) = 200$ s; median 137 s.
\textbf{4.} 15.6\%.
\textbf{5.} One: the queue drifts down at 0.1 order a second (99.6\% of paths within 20 minutes).
\textbf{6.} 73.2\%.
\textbf{7.} 73.4\%.
\textbf{8.} 99.0\% with 40 orders on the ask; with 80, the ask practically never empties first.
\textbf{9.} 99.3\%.
\textbf{10.} 95.4\% up with 80 against 20; 50\% with equal queues.
\textbf{11.} 10.8\%.
\textbf{12.} 156.6 seconds.
\textbf{13.} Departures would grow with the queue, pulling long queues down and short ones up: a mean-reverting size
instead of a constant drift, and fewer extreme queues.
\textbf{14.} Bursts make some queues empty much faster and others last longer: fatter tails on both sides, and a fill
probability that depends on the state of the flow.
\textbf{15.} The dependence between the sides: a large buy empties the ask and brings in new bids, and liquidity
providers react to imbalance (the queue-reactive models of One Quant Book~10).
\textbf{16.} The difference in fill probability between the front and the back, $100\%$ against $73\%$ here, times the
value of a fill (half the spread less the expected adverse move).
\textbf{17.} That the bid was hit 80 times before the ask emptied: sellers were aggressive, and the price is more likely
to fall next than the unconditional numbers say.
\textbf{18.} Only if half the spread exceeds the expected adverse move after a fill (the mark-out, One Quant Book~2,
chapter~15); the model gives the fill probability, not the answer.
\textbf{19.} Named result: \emph{the back of the queue}: with 80 orders ahead and 20 on the ask, the order fills before
the ask empties with probability 73\% (73.2\% exact, 73.4\% simulated), after 82 seconds on average; with the ask
twice as long, 99\%.
\textbf{20.} The race between the orders ahead leaving and the other side's queue emptying.
\end{solution}

\begin{solution}{iq:qm:markov-chains-and-queues:1}
As an \omterm{def:qm:markov-chains-and-queues:mm1}{M/M/1 queue} with $\rho = 0.9$: 9 requests in the system on average (8.1 waiting), each spending $1/(10\,000 - 9\,000)$
s $= 1$ ms, ten times the service time.
\iqlookfor{$\rho/(1-\rho)$ and Little's law, and the blow-up near saturation.}
\end{solution}

\begin{solution}{iq:qm:markov-chains-and-queues:2}
The matrix of jump rates $q_{ij}$, with diagonal minus the row sums; $P(t) = e^{t\mathcal L}$. The \omterm{def:qm:stochastic-differential-equations:kolmogorov}{stationary distribution}
solves $\pi\mathcal L = 0$ with $\sum\pi_i = 1$; for a birth--death chain, \omterm{def:qm:markov-chains-and-queues:balance}{detailed balance} gives it in closed form.
\iqlookfor{rates, holding times, and $\pi\mathcal L = 0$.}
\end{solution}

\begin{solution}{iq:qm:markov-chains-and-queues:3}
Up: the thin ask empties sooner than the heavy bid, and the first queue to empty moves the price. With the
chapter's rates, 95\%.
\iqlookfor{queue depletion as a race.}
\end{solution}

\begin{solution}{iq:qm:markov-chains-and-queues:4}
$3/10$, by optional stopping; the expected number of steps is $3 \times 7 = 21$, from $W_n^2 - n$.
\iqlookfor{gambler's ruin and $i(n - i)$.}
\end{solution}

\begin{solution}{iq:qm:markov-chains-and-queues:5}
The front fills first and more often, before the price moves away, and the back fills mainly when the queue is
being eaten by informed flow. Value it as the difference in fill probability times the edge per fill, with fill
probabilities from a queue model or from the firm's own fill data by queue position.
\iqlookfor{fill probability by position, and adverse selection.}
\end{solution}

\begin{solution}{iq:qm:markov-chains-and-queues:6}
Uniformisation: pick $L$ at least every exit rate, run the jump matrix $I + \mathcal L/L$ on a vector (sparse
matrix--vector products) and weight the iterates by Poisson probabilities. Every term is nonnegative, so there is
no cancellation, unlike a Taylor series of $e^{t\mathcal L}$.
\iqlookfor{uniformisation, sparsity, and positivity.}
\end{solution}
